\section{Objetivos del proyecto}

\subsection{Objetivo general}

Diseñar e implementar un sistema de automatización distribuido basado en el protocolo \textit{CAN BUS} para integrar nuevas unidades de control en un vehículo antiguo, ampliando sus funcionalidades hacia áreas de confort, multimedia y conectividad avanzada, incluyendo interacción con una aplicación móvil y control por voz mediante inteligencia artificial.

\subsection{Objetivos específicos}

\begin{itemize}
    \item Analizar el sistema electrónico original del vehículo y su nivel de conectividad mediante el puerto OBD-II.
    \item Diseñar e implementar una red CAN adicional que permita la comunicación entre nodos distribuidos.
    \item Desarrollar una \textbf{unidad de confort} encargada de gestionar funciones como climatización, cierre centralizado, control de ventanas y ajustes de iluminación interior.
    \item Desarrollar una \textbf{unidad multimedia} enfocada en la gestión de audio, interfaces de usuario y notificaciones.
    \item Implementar controladores basados en el chip \textbf{MCP2515} y microcontroladores Atmega para asegurar la comunicación estandarizada dentro de la red CAN.
    \item Incorporar una plataforma de procesamiento avanzado (p. ej., Raspberry Pi) para la gestión de interfaces multimedia, comunicación externa y servicios inteligentes.
    \item Diseñar e integrar una \textbf{aplicación móvil} que permita el monitoreo y control remoto de funciones del vehículo a través de la red CAN.
    \item Desarrollar un sistema de \textbf{control por voz basado en inteligencia artificial}, capaz de ejecutar comandos del usuario para gestionar funciones de confort y multimedia.
    \item Validar experimentalmente el funcionamiento del sistema mediante pruebas de integración, conectividad móvil y control por voz en condiciones reales.
    \item Documentar el proceso de desarrollo y generar una bitácora de experimentos que permita futuras mejoras y replicabilidad del sistema.
\end{itemize}
